%% %%%%%%%%%%%%%%%%%%%%%%%%%%%%%%%%%%%%%%%%%%%%%%%%
%% Problem Set/Assignment Template to be used by the
%% Food and Resource Economics Department - IFAS
%% University of Florida's graduates.
%% %%%%%%%%%%%%%%%%%%%%%%%%%%%%%%%%%%%%%%%%%%%%%%%%
%% Version 1.0 - November 2019
%% %%%%%%%%%%%%%%%%%%%%%%%%%%%%%%%%%%%%%%%%%%%%%%%%
%% Ariel Soto-Caro
%%  - asotocaro@ufl.edu
%%  - arielsotocaro@gmail.com
%% %%%%%%%%%%%%%%%%%%%%%%%%%%%%%%%%%%%%%%%%%%%%%%%%

\documentclass[12pt]{article}
\usepackage{design_ASC}

\setlength\parindent{0pt} %% Do not touch this

%% -----------------------------
%% TITLE
%% -----------------------------
\title{Problem Set \#1} %% Assignment Title

\author{John Nash\\ %% Student name
AEB6933 - Fantastic Econometric Course\\ %% Code and course name
\textsc{University of Florida}
}

\date{\today} %% Change "\today" by another date manually
%% -----------------------------
%% -----------------------------

%% %%%%%%%%%%%%%%%%%%%%%%%%%
\begin{document}
\setlength{\droptitle}{-5em}    
%% %%%%%%%%%%%%%%%%%%%%%%%%%
\maketitle

% --------------------------
% Start here
% --------------------------

% %%%%%%%%%%%%%%%%%%%
\section*{Question 1}
% %%%%%%%%%%%%%%%%%%%
{\bfseries You can write the questions here \ldots}

And your answers here \ldots

% %%%%%%%%%%%%%%%%%%%
\subsection*{Question 1.a}
% %%%%%%%%%%%%%%%%%%%
{\bfseries Also, you can write your sub-questions into subsection levels.}

And again, your answers here \ldots

% %%%%%%%%%%%%%%%%%%%
\section*{Mathematical notation}
% %%%%%%%%%%%%%%%%%%%

This template offers a set of customized mathematical symbols commonly used in FRED courses.

\begin{multicols}{2}
You can invoke the following symbols:
\begin{Verbatim}[frame=single, fontsize=\footnotesize]
\E(\cdot)             % Expectation
\V(\cdot)             % Variance
\Var(\cdot)           % Variance
\Cov(\cdot)           % Covariance
\Corr(\cdot)          % Correlation
\tr(\cdot)            % Trace
\rank(\cdot)          % Rank
\N(0,1)               % Normal
\op(\cdot)            % "little o" ope
\Op(\cdot)            % "big o" operator
\R^+                  % Real positive num
\La(y|\beta,\sigma^2) % Lagrange function
\end{Verbatim}
\footnotesize{
\begin{itemize}  \setlength\itemsep{0em}
    \item $\E(\cdot)$
    \item $\V(\cdot)$
    \item $\Var(\cdot)$
    \item $\Cov(\cdot)$
    \item $\Corr(\cdot)$
    \item $\tr(\cdot)$
    \item $\rank(\cdot)$
    \item $\N(0,1)$
    \item $\op(\cdot)$
    \item $\Op(\cdot)$
    \item $\R^{+}$
    \item $\La(y|\beta,\sigma^2)$
\end{itemize}}
\end{multicols}

Remember, these symbols have to be implemented into equation  environments, that means, with dollar symbols \$\ldots\$, or \texttt{equation}-environments.

\newpage 

There is a set of different accents you may use:
\begin{Verbatim}[frame=single, fontsize=\footnotesize]
 $\hat \beta$, $\widehat \beta$, $\tilde \alpha$, $\widetilde \alpha$, $\bar y$, 
 $\varepsilon \sim \chi^2_t$
\end{Verbatim}
Which reproduce the following: $\hat \beta$, $\widehat \beta$, $\tilde \alpha$, $\widetilde \alpha$, $\bar y$, $\varepsilon \sim \chi^2_t$.

These are a couple of examples:
\begin{equation}
    \Avar (b) = \frac{\sigma^2}{n} Q^{-1} \plim \left( \frac{1}{n}X'QX \right)Q^{-1}
\end{equation}

\begin{equation}
    \sqrt{n}(b-\beta) \stackrel{d}{\longrightarrow} \N \left[0,\frac{\sigma^2}{n} Q^{-1} \plim \left( \frac{1}{n}X'\Omega X\right)Q^{-1}\right]
\end{equation}



% %%%%%%%%%%%%%%%%%%%
\section*{Code and scripts}
% %%%%%%%%%%%%%%%%%%%

If you want to print the raw outcomes from any software, it is recommended the \texttt{Verbatim}-environment :
%% Verbatim -----------------
\begin{Verbatim}[frame = single, fontsize = \footnotesize]
Parameters    Estimates     Std. err.  Est./s.e.  Prob.    Gradient
------------------------------------------------------------------
P01             -0.5553        0.1446   -3.840   0.0001      0.0002
P02             -0.2255        0.1106   -2.039   0.0414     -0.0004
P03             -0.8588        0.2378   -3.612   0.0003     -0.0003
\end{Verbatim}
%% Verbatim -----------------
The same environment can be used for print a chunk of code (this is GAUSS by the way):
%% Verbatim -----------------
\begin{Verbatim}[frame = single, fontsize = \footnotesize]
ev = ev1 ~ ev2;
ev1 = sumc(ev[.,nest1]');
ev2 = sumc(ev[.,nest2]');
num = (ev1 .^ (k[1]-1)).*sumc(depm[.,nest1]') + (ev2 .^ (k[2]-1)).*sumc(depm[.,nest2]');
p = sumc((ev .* depm)').* num ./ ((ev1.^k[1])+(ev2.^k[2]));
\end{Verbatim}
%% Verbatim -----------------
If you want to print Stata-code, you can use the environment \texttt{lstlisting}, setting the \texttt{style} in \texttt{stata-editor}. This will highlight the Stata commands like the following example:
%% Stata-code -----------------
\begin{lstlisting}[language = Stata, style = stata-editor]
clogit depvar var1 var2 var3, group(id)
\end{lstlisting}
%% Stata-code -----------------
The same \texttt{lstlisting}-package can be used to print \textbf{R}-code, setting the only argument (\texttt{language}) in \texttt{R}.
%% R-code -----------------
\begin{lstlisting}[language = R]
Quad.reg <- function(z) {
  grad <- cbind(z[4]+z[11]*m[1]+z[15]*m[2]+z[25]*m[3]+z[19]*m[4]+z[20]*m[5]+z[21]*m[6],
                z[5]+z[12]*m[1]+z[16]*m[2]+z[19]*m[3]+z[26]*m[4]+z[22]*m[5]+z[23]*m[6],
                z[7]+z[14]*m[1]+z[18]*m[2]+z[21]*m[3]+z[23]*m[4]+z[22]*m[5]+z[28]*m[6])
  test1 <- min(grad) # have to be >0
  test2 <- max(eigen(hess)$values) # have to be < 0
  if(Mode == 1) {
    lst <- list(Grad = grad, Hess = hess, Test1 = test1, Test2 = test2)
    return(lst)
  } else {
    return(cbind(test1,test2))
  }
}
\end{lstlisting}
%% R-code -----------------




\end{document}